\chapter{Access, Identity and Auditability}
\section*{Introduction}
\addcontentsline{toc}{section}{Introduction}
Access is the first technical control in \projectname{} because every later module depends on knowing who is using the platform and which organisational scope they are allowed to see. This chapter therefore covers three connected concerns: authentication, authorization, and auditability. Authentication establishes a verified session, authorization turns that session into page and API permissions, and auditability preserves a record of sensitive partner lifecycle decisions.
\section{Sprint and module objective}
This sprint turns the platform from a set of public and private screens into a governed workspace. Its scope is deliberately foundational: before partner CRM, HR workflows, dashboards, reports, or AI tools can be trusted, the system must authenticate the user, route that user to the correct workspace, and enforce a consistent permission model.
\subsection{Functional objective}
The objective of this module is to provide a secure entry point for all IntelliConnect users while keeping public pages, internal employee operations and partner self-service separate.
\begin{itemize}
  \item provide a public entry surface for visitors without exposing protected data;
  \item authenticate employees and partners with email and password credentials;
  \item issue a server-verifiable JWT session stored in an httpOnly cookie;
  \item extract the session on the server before rendering protected layouts or processing protected API requests;
  \item route employees and partners toward the correct portal after sign-in;
  \item enforce role-based permissions for Admin, Manager, Commercial, Analyst, HR and Partner users;
  \item preserve auditable records for partner status changes and other sensitive actions.
\end{itemize}
\subsection{Backlog for access, identity and auditability}
\begin{table}[H]
  \centering
  \caption{Backlog of the access, identity and auditability module}
  \begin{tabularx}{\textwidth}{>{\bfseries\raggedright\arraybackslash}p{1.2cm}>{\raggedright\arraybackslash}p{3.1cm}X>{\raggedright\arraybackslash}p{2.1cm}}
    \toprule
    ID & Theme & User story & Status \\
    \midrule
    AI-1 & Public access & As a visitor, I can consult public pages without authentication while protected information remains inaccessible. & Completed \\
    AI-2 & Sign-in & As an authorized user, I can sign in with an email and password through a unified access page. & Completed \\
    AI-3 & Session security & As the system, I store the authenticated session in an httpOnly cookie and verify it server-side with \sourcecode{jose}. & Completed \\
    AI-4 & Password security & As the system, I store only hashed passwords using \sourcecode{bcryptjs} and never persist plain-text credentials. & Completed \\
    AI-5 & Portal routing & As an authenticated user, I am redirected to the employee or partner portal according to my account type and role. & Completed \\
    AI-6 & RBAC & As an administrator, I can rely on role-based permissions to restrict sensitive screens and operations. & Completed \\
    AI-7 & API protection & As the backend, I reject unauthenticated, unauthorized or invalid requests with structured API errors. & Completed \\
    AI-8 & Auditability & As a governance user, I can trace partner status changes with actor, previous status, new status, reason and timestamp. & Completed \\
    \bottomrule
  \end{tabularx}
\end{table}
\section{Implementation analysis}
The implementation deliberately keeps authentication logic server-side. Client pages can display forms and loading states, but they do not decide whether a user is trusted, which role is active, or which partner records can be returned. Those decisions are made by route handlers, layouts, and service functions that read the signed session and enforce the same boundary before any protected data leaves the server.
\subsection{JWT sessions with httpOnly cookies}
User sessions are represented by signed JSON Web Tokens issued after a successful credential check. The token is stored in an httpOnly cookie rather than browser storage, which reduces exposure to client-side script access and keeps session handling compatible with server-rendered layouts. The cookie carries only the session claims needed by the application: the user identifier, account type, role, and partner scope when the authenticated user belongs to an external partner organisation.

This design gives the platform a single portable session format for employee and partner flows. A dashboard request, a partner portal request, and a protected API call all start from the same primitive: read the cookie, verify the signature, recover the claims, then decide whether the requested operation is allowed.
\subsection{Session extraction and server-side trust boundary}
The trust boundary is crossed only after verification on the server. In the application layer, session helper functions extract the cookie, validate the token with the configured secret, and normalize the result into a small authenticated-user object. If extraction fails, the protected layout redirects the user to sign in, while API routes return a structured \sourcecode{401} response.

Keeping this extraction centralized avoids duplicating authentication checks across pages. It also prevents a common class of inconsistencies where a route appears protected in the interface but still exposes data through a direct API call. In \projectname{}, protected rendering and protected data access share the same session prerequisite.
\subsection{Password hashing with bcryptjs}
Passwords are never stored in clear text. During account creation or demo seeding, the backend hashes the password with \sourcecode{bcryptjs}; during sign-in, the submitted password is compared against that stored hash. The database therefore contains a verifier, not the original secret.

The practical consequence is that authentication remains recoverable and maintainable without weakening the storage layer. Administrators can reset or seed demo credentials when needed, but the platform does not need to read or display a user's existing password. Failed sign-in attempts return generic messages so the interface does not reveal whether the email address or the password was the invalid part.
\subsection{Route protection and role-aware rendering}
Route protection combines the authenticated session with role-aware rendering. Employee pages live under the internal dashboard area and are only available to employee roles such as Admin, Manager, Commercial, Analyst, and HR. Partner pages live under the partner portal and are scoped to the authenticated partner organisation. The user interface follows this split by showing only the navigation and actions relevant to the current role.

The same rule is repeated at the API boundary. A hidden button is not treated as a security control: protected route handlers verify the session and the role before executing the operation. This is especially important for partner records, uploaded documents, status transitions, reports, and AI-assisted analysis, where the user may see different data depending on their role or partner scope.
\subsection{Structured API errors}
\begin{verbatim}
{ "error": "Validation failed", "details": [...] }
\end{verbatim}
The status code reflects the failure category: \sourcecode{401} for missing or invalid authentication, \sourcecode{403} for insufficient permissions, \sourcecode{400} for invalid input and \sourcecode{500} for unexpected server failures.
This error format keeps the frontend predictable. Forms can attach validation details to fields, dashboards can show a recoverable error state, and the agent interface can explain that a protected tool failed without exposing stack traces, SQL details, or secrets. The implementation therefore treats error shape as part of the access contract rather than an incidental response format.
\section{Role-based access control}
Role-based access control translates the authenticated identity into concrete application permissions. The module separates three questions that are often mixed together: whether the user is signed in, whether the user belongs to the employee or partner side of the platform, and whether the user's role authorizes the requested action. This separation keeps the public website, employee workspace, and partner self-service portal from leaking into one another.
\subsection{Access-control flow}
\begin{figure}[htbp]
  \centering
  \dgram{access-flow}
  \caption{Access-control decision applied to every protected request.}
  \label{fig:access-flow}
\end{figure}
The flow is intentionally evaluated on every protected request, not only at sign-in time. A valid cookie proves that a user authenticated previously; it does not by itself prove that the user can access a particular page, partner record, report, or administrative action. The role and ownership check is therefore repeated where the data is requested.
\subsection{RBAC role matrix}
The role matrix below documents the operational permissions used by the platform. It is written from the business capability point of view rather than from file paths, because the same capability can appear in multiple screens or APIs. Partner users are marked separately because their access is limited to their own organisation even when a feature name resembles an employee feature.
\begin{table}[H]
  \centering
  \caption{RBAC matrix for IntelliConnect roles}
  \begin{tabularx}{\textwidth}{>{\bfseries\raggedright\arraybackslash}p{3.2cm}cccccc}
    \toprule
    Capability & Admin & Manager & Commercial & Analyst & HR & Partner \\
    \midrule
    Access employee dashboard & Yes & Yes & Yes & Yes & Yes & No \\
    Access partner portal & No & No & No & No & No & Yes \\
    Manage user accounts and roles & Yes & No & No & No & No & No \\
    View partner portfolio & Yes & Yes & Yes & Yes & Yes & Limited \\
    Create and update partner records & Yes & Yes & Yes & No & No & Limited \\
    Review partnership requests & Yes & Yes & Yes & No & No & No \\
    View analytics and BI dashboards & Yes & Yes & Yes & Yes & Yes & Limited \\
    Manage HR workflows & Yes & No & No & No & Yes & No \\
    Upload or manage partner documents & Yes & Yes & Yes & No & No & Own scope \\
    View audit and status history & Yes & Yes & Yes & Yes & Yes & Own scope \\
    Use AI decision-support tools & Yes & Yes & Yes & Yes & Yes & No \\
    \bottomrule
  \end{tabularx}
\end{table}
\subsection{Employee and partner portal separation}
\begin{itemize}
  \item The employee portal, exposed under \sourcecode{/dashboard}, centralizes internal operations: partner records, contacts, offers, events, projects, documents, status history, reports, BI and AI-assisted analysis.
  \item The partner portal, exposed under \sourcecode{/partner}, provides a self-service workspace for external organizations: profile, contacts, notifications, documents, offers, events, statistics and category-specific collaboration features.
\end{itemize}
This separation matters for both security and usability. Employees need a broad operational view across the portfolio, while partners need a narrow workspace that exposes their own documents, offers, notifications, and collaboration data without revealing internal review notes or other organisations' records.
\subsection{Public routes and protected routes}
\begin{table}[H]
  \centering
  \caption{Public and protected route categories}
  \begin{tabularx}{\textwidth}{>{\bfseries\raggedright\arraybackslash}p{3.6cm}
                               >{\raggedright\arraybackslash}X}
    \toprule
    Category & Access rule \\
    \midrule
    Public routes & Accessible without authentication. The public website and legal pages serve visitors without exposing any protected operational data. \\
    Authentication route & Available to unauthenticated users. After a successful sign-in, the server routes the user to the correct portal based on their identity type. \\
    Employee protected routes & Require a valid employee session and an authorized employee role. Access to each page is scoped to the user's role in the RBAC matrix. \\
    Partner protected routes & Require a valid partner session. Every request is restricted to data belonging to the authenticated partner organization. \\
    Protected API routes & Require session extraction, role verification, input validation, and return structured errors. Raw stack traces and database details are never exposed. \\
    \bottomrule
  \end{tabularx}
\end{table}
The route categories make the implementation easier to audit. Public pages are allowed to be anonymous by design; everything else must justify its access path through a session and, when relevant, a role or ownership check. This gives reviewers a simple checklist when adding new routes: decide the category first, then apply the matching guard before returning data.
\section{Auditability and traceability}
Auditability completes the access model. It is not enough to know that a user was allowed to perform an action; for partnership management, the platform also needs to preserve who changed a business state, when it changed, what the previous value was, and why the change was made. These records support internal review, traceability during demonstrations, and later governance work if the platform is industrialized.
\subsection{Partner status history}
\begin{table}[H]
  \centering
  \caption{Information retained for partner status history}
  \begin{tabularx}{\textwidth}{>{\bfseries\raggedright\arraybackslash}p{3.2cm}X}
    \toprule
    Field & Purpose \\
    \midrule
    Actor & Identifies the authenticated user who performed the transition. \\
    Previous status & Preserves the state before the operation, enabling reconstruction of the lifecycle path. \\
    New status & Records the resulting business state after the operation. \\
    Reason & Captures the justification entered by the user or produced by the workflow. \\
    Timestamp & Establishes when the transition occurred for chronological and legal traceability. \\
    \bottomrule
  \end{tabularx}
\end{table}
The retained fields are deliberately business-readable. They allow a reviewer to reconstruct a partner's lifecycle without reading raw database logs: the actor identifies responsibility, the previous and new statuses show the transition, the reason explains the decision, and the timestamp orders the event in the broader partnership timeline.
\subsection{Action logs and status transitions}
Action logs complement status history by capturing sensitive operations around the partner lifecycle. Typical examples include accepting a partnership request, rejecting an incomplete application, updating a partner's status, or uploading a document that supports a decision. These events are stored as application records, which makes them easier for managers to review than infrastructure logs.

The objective is not to record every low-level click. The platform focuses on actions that change business state, affect access, or support a decision. This keeps the audit trail useful: dense enough to explain the lifecycle, but not so noisy that important transitions disappear inside routine navigation events.
\subsection{Security and legal traceability}
Security traceability also supports legal and organisational accountability. Partner data includes contacts, documents, offers, and project information, so the platform must be able to explain why a user saw or changed a record. By combining role checks, ownership checks, and status history, \projectname{} can show that protected operations followed an explicit access rule.
\section{Implementation}
The implementation artefacts below connect the access model to the visible product. They show how the authentication use cases, sign-in surface, and role-routing sequence appear in the delivered platform. The screenshots and diagrams are not separate features; they are the concrete entry points through which the previous security rules are exercised.
\subsection{Access use cases}
\begin{figure}[htbp]
  \centering
  \dgram{usecase-auth}
  \caption{Use cases of the access and identity module.}
\end{figure}
The use-case diagram summarizes the module from the user's point of view. Public visitors can access only public information, employees authenticate into the internal workspace, partners authenticate into their own portal, and administrators retain the extra responsibility of managing roles and sensitive account operations.
\subsection{Sign-in page}
\begin{figure}[htbp]
  \centering
  \shot{signin}
  \caption{Sign-in page used as the controlled entry point to IntelliConnect.}
\end{figure}
The sign-in page is the controlled entry point for both employee and partner accounts. It collects credentials once, delegates verification to the server, and then sends the authenticated user to the correct workspace. Keeping the entry point unified reduces confusion while the post-login routing still preserves the employee--partner separation.
\subsection{Authentication and role-routing sequence}
\begin{figure}[htbp]
  \centering
  \dgram{sequence-auth}
  \caption{JWT creation, session extraction, and role-based portal routing.}
\end{figure}
The sequence diagram shows the same control path in chronological order. The user submits credentials, the server verifies the password hash, creates a signed JWT, stores it in the httpOnly cookie, and resolves the target portal from the authenticated role. Subsequent requests do not repeat password verification; they verify the session and apply the role and ownership rules described earlier.
\section*{Conclusion}
\addcontentsline{toc}{section}{Conclusion}
This chapter establishes the access, identity and auditability foundation of \projectname{}. The module combines JWT sessions stored in httpOnly cookies, server-side verification with \sourcecode{jose}, password hashing with \sourcecode{bcryptjs}, protected routing, structured API errors and a role matrix adapted to employee and partner responsibilities.
